\chapter{Debate de artículos de RRHH en los equipos}

En esta dinámica trabajamos con distintos artículos que abordaban temas muy diversos, desde la negociación laboral hasta la importancia del \emph{networking} o la definición de una estrategia vital. Lo que más me llamó la atención fue que, aunque cada texto trataba un aspecto diferente, todos coincidían en un mensaje de fondo: \textbf{nuestra carrera profesional no puede basarse solo en el corto plazo ni en criterios exclusivamente económicos, sino que debe estar guiada por valores, propósito y relaciones de calidad}. 

Uno de los debates más enriquecedores fue el del artículo \emph{Negotiating as a Woman of Color} \cite{hbr-woc-negotiation}. En clase se comentó cómo los prejuicios pueden convertirse en un obstáculo, pero también en una oportunidad si se saben gestionar con inteligencia. Esto me recordó al libro \emph{The Broken Rung} \cite{broken-rung}, donde se analiza cómo mujeres y personas racializadas suelen tener muchos mentores que dan consejos, pero menos patrocinadores que crean oportunidades reales. Más allá de la anécdota, me hizo reflexionar sobre la importancia de que las organizaciones cuenten con políticas que aseguren la igualdad de oportunidades y que reconozcan el valor de la diversidad como motor de crecimiento. Esta idea conecta con lo trabajado en clase sobre la necesidad de promover la igualdad y la diversidad como elementos clave en la gestión de personas.  

El tema del \textbf{networking} apareció en el artículo \emph{Learn to Love Networking} \cite{hbr-networking}. La profesora insistió en que las redes profesionales no deben basarse únicamente en pedir, sino en compartir y ser generoso, resumiéndolo con la frase “nunca comas solo”. Me quedo con la idea de que el \emph{networking} no es un acto interesado, sino una forma de crear vínculos de confianza que facilitan el trabajo en equipo y la innovación. También pensé en cómo, tal y como vimos en clase, las relaciones en red y las comunidades de práctica son cada vez más relevantes para compartir experiencias y conocimiento, lo que a la larga se traduce en un entorno más creativo y colaborativo.

La metáfora del \textbf{marshmallow}, presentada en el artículo \emph{How to Choose a Fulfilling Career} \cite{atlantic-fulfilling-career}, ilustraba la importancia de la paciencia y la visión a largo plazo. El profesor lo relacionó con el hecho de que el propio trabajo debe ser fuente de satisfacción y no solo un medio para obtener un sueldo o un cargo. Para mí fue revelador, porque a menudo pienso en la meta final y olvido que también debo disfrutar del camino. En este punto recordé otro consejo leído: no hay que tener miedo a empezar de nuevo, reinventarse o incluso hacer movimientos laterales en la carrera. Esto lo conecto con lo visto en clase sobre la empleabilidad: la carrera profesional no siempre es lineal, sino que requiere flexibilidad y capacidad de adaptación para mantenerse vigente en un entorno cambiante.  

La reflexión sobre las \textbf{80.000 horas de trabajo}, propuesta en \emph{This Is Your Most Important Decision} \cite{80000hours}, fue otro de los puntos que me impactó. Tomar conciencia de que dedicaremos tanto tiempo a nuestra vida profesional es un buen recordatorio de que debemos invertirlo en algo que tenga sentido y contribuya a la sociedad. También pensé que, para que esto sea posible, las empresas necesitan diseñar entornos que motiven, retengan y desarrollen a las personas. Esto enlaza con lo trabajado en clase sobre los objetivos de la gestión de personas: atraer, retener, motivar y desarrollar. Un empleo con propósito no solo beneficia al trabajador, sino también a la organización, que obtiene compromiso y mejores resultados cuando las personas sienten que lo que hacen tiene valor.

Finalmente, artículos como \emph{How Will You Measure Your Life?} \cite{christensen-measure-life} y \emph{Negotiating Your Next Job} \cite{hbr-negotiating-next-job} reforzaron esta visión integral. Aprendí que no se trata de aceptar cualquier oferta pensando solo en el presente, sino de diseñar una estrategia coherente con mis objetivos y principios; y que, al negociar, conviene mirar también el rol, las oportunidades de desarrollo y el encaje con mis metas. Esto me hizo pensar en lo trabajado en clase sobre el contrato psicológico: ese vínculo no escrito entre empleado y organización que abarca mucho más que el salario, e incluye confianza, expectativas de crecimiento y cultura organizacional.  

En conclusión, esta dinámica me ayudó a mirar la carrera profesional desde una perspectiva más amplia y personal. Entendí que el éxito no se mide solo en cifras o ascensos, sino en la coherencia entre mis valores, el propósito que guíe mis decisiones y las relaciones que construya a lo largo del camino, sin perder de vista esos momentos clave que pueden marcar la diferencia en una trayectoria. 
